\begin{tikzpicture}[scale=1]
    % 定义圆盘的半径
    \def\R{5}
    
    % --- 宏：计算并绘制分数节点 ---
    \newcommand{\fnode}[2]{
        % 角度变换映射：2 * arctan(y/x)
        \pgfmathsetmacro{\ang}{2*atan2(#1, #2)}
        \node[font=\scriptsize, inner sep=1pt] at (\ang:\R+0.45) {$#1/#2$};
    }
    
    % --- 宏：计算并绘制两条测地线之间的双曲圆弧 ---
    \newcommand{\drawgeodesic}[4]{
        \pgfmathsetmacro{\angA}{2*atan2(#1, #2)}
        \pgfmathsetmacro{\angB}{2*atan2(#3, #4)}
        \pgfmathsetmacro{\diff}{abs(\angA - \angB)}
        
        % 判断是否为穿过圆心的直线（即角度差为 180 度）
        \pgfmathparse{\diff < 179.9 && \diff > 0.1 ? 1 : 0}
        \ifnum\pgfmathresult=1
            % 计算正交圆的圆心和半径
            \pgfmathsetmacro{\midang}{(\angA + \angB)/2}
            \pgfmathsetmacro{\halfang}{(\angA - \angB)/2}
            \pgfmathsetmacro{\rad}{\R * abs(tan(\halfang))}
            \pgfmathsetmacro{\cx}{\R * cos(\midang)/cos(\halfang)}
            \pgfmathsetmacro{\cy}{\R * sin(\midang)/cos(\halfang)}
            \draw[semithick] (\cx,\cy) circle (\rad);
        \else
            % 如果是直线，直接连接两个点
            \draw[thick] (\R*cos \angA, \R*sin \angA) -- (\R*cos \angB, \R*sin \angB);
        \fi
    }
    
    % --- 宏：递归生成 Farey 序列并绘制网格 ---
    \newcommand{\farey}[5]{
        % #1/#2 和 #3/#4 是端点，#5 是递归深度
        \ifnum#5>0
            % 计算 Farey 中间分数
            \pgfmathtruncatemacro{\p}{#1+#3}
            \pgfmathtruncatemacro{\q}{#2+#4}
            
            % 绘制连线
            \drawgeodesic{#1}{#2}{\p}{\q}
            \drawgeodesic{\p}{\q}{#3}{#4}
            \fnode{\p}{\q} % 绘制节点数字
            
            % 向下递归
            \pgfmathtruncatemacro{\newdepth}{#5-1}
            \farey{#1}{#2}{\p}{\q}{\newdepth}
            \farey{\p}{\q}{#3}{#4}{\newdepth}
        \fi
    }
    
    % --- 绘图主体 ---
    \begin{scope}
        % 限制所有线条在基本圆盘内部
        \clip (0,0) circle (\R);
        
        % 绘制作为基础框架的 4 条基准线
        \drawgeodesic{1}{0}{0}{1}
        \drawgeodesic{0}{1}{1}{1}
        \drawgeodesic{1}{1}{1}{0}
        \drawgeodesic{-1}{0}{-1}{1} % 注意: -1/0 对应几何上的 1/0
        \drawgeodesic{-1}{1}{0}{1}
        
        % 递归深度设置为 4，刚好对应你原图中的分数密度
        \farey{0}{1}{1}{1}{4}
        \farey{1}{1}{1}{0}{4}
        \farey{-1}{0}{-1}{1}{4}
        \farey{-1}{1}{0}{1}{4}
    \end{scope}
    
    % 绘制最外侧的边界圆（加粗）
    \draw[thick] (0,0) circle (\R);
    
    % 绘制 4 个最基础的极点节点
    \fnode{0}{1}
    \fnode{1}{1}
    \fnode{1}{0}
    \fnode{-1}{1}
    \fnode{-3}{1}
    \fnode{-1}{2}
    
\end{tikzpicture}